\newpage
\section{Theoretische Analyse}

\subsection{Theoretischer Hintergrund}

\subsubsection{Effizienzmarkthypothese (EMH) und Grenzen der Prognostizierbarkeit}
\paragraph{Grundannahme}
Ein zentrales Grundprinzip moderner Kapitalmarktmodelle ist, dass Finanzmarktpreise alle öffentlich verfügbaren Informationen widerspiegeln sollten, ist eine Annahme, auf die Effizienzmarkthypothese (EMH) aufbaut. Danach sollten insbesondere neue, unerwartete Informationen zu wirtschaftlichen Fundamentaldaten, wie Konjunkturdaten oder geldpolitischen Entscheidungen, die treibenden Kräfte hinter Preisänderungen an Finanzmärkten sein.

In der empirischen Forschung zeigte sich jedoch, dass klassische Untersuchungen häufig nur einen geringen Anteil der beobachteten Marktbewegungen auf solche Nachrichten zurückführen konnten, was den Eindruck erweckte, Finanzmarktpreise seien weitgehend von den Fundamentaldaten entkoppelt. Die Bundesbank-Studie erweitert die Datenbasis: Sie nutzt eine umfassende, mit Zeitstempeln versehene Ereignisdatenbank, die sowohl planbare makroökonomische Veröffentlichungen als auch unvorhergesehene Ereignisse wie Rating-Herabstufungen, Wahlen, Kriege oder Naturkatastrophen umfasst, und wertet deren Zusammenhang mit hochfrequenten Preisänderungen aus.

Die Ergebnisse zeigen, dass etwa die Hälfte aller hochfrequenten Marktbewegungen rund um eindeutig identifizierbare Nachrichten stattfindet und sich somit in diesem Sinne durch Informationen erklären lässt. Damit fällt der Anteil der erklärten Preisbewegungen deutlich höher aus als in vielen früheren Studien. Das unterstützt die Idee, dass Finanzmarktpreise nicht vollständig von Fundamentaldaten losgelöst sind, sondern dass sich zumindest ein bedeutender Teil der Volatilität auf reale Informationsübermittlungen zurückführen lässt. Allerdings bleibt offen, wie Nachrichten zu Preisbewegungen führen: Sie könnten entweder direkt Erwartungen über zukünftige Zahlungsströme oder Diskontfaktoren ändern oder aber über indirekte Marktreaktionen und Transaktionen Preisbewegungen auslösen.


\subsubsection{Marktpsychologie und Informationsasymmetrien}
\paragraph{Problemstellung}
Ein wichtiges ökonomisches Problem, das eng mit der Bewertung und Prognose von Aktienbewegungen zusammenhängt, ist die Informationsasymmetrie. Informationsasymmetrie bedeutet, dass nicht alle Marktteilnehmer gleich viele oder gleich gute Informationen besitzen. Manche wissen mehr über bestimmte Risiken oder Chancen als andere, was zu einem ungleichen Informationsstand führt.

Ein klassisches Beispiel dafür ist das Konzept der „Adverse Selection“. Dieser Begriff beschreibt eine Situation, bei der eine Marktseite bessere Informationen hat als die andere. Dadurch kommt es zu einem Marktversagen, weil die schlechter informierten Akteure Entscheidungen treffen müssen, ohne wirklich zu wissen, was sie bekommen. In dem berühmten Modell von George Akerlof ("The Market of Lemons") können Käufer auf einem Gebrauchtwagenmarkt die Qualität der Autos nicht erkennen, während die Verkäufer genau Bescheid wissen. Das führt dazu, dass nur noch Autos niedriger Qualität angeboten werden, weil gut erhaltene Fahrzeuge vom Markt verschwinden.

Übertragen wir diese Idee auf den Aktienmarkt: Wenn einige Investoren insiderähnliche Informationen besitzen oder genauer einschätzen können, wie es um ein Unternehmen steht, haben sie einen Vorteil gegenüber anderen. Die Folge ist, dass weniger informierte Anleger falsche Erwartungen über die zukünftige Kursentwicklung haben können. Gerade bei der Prognose von Aktienkursen beeinflusst diese unterschiedliche Informationsverteilung die Marktpreise und die Marktpsychologie stark. Anleger, die glauben, sie hätten verlässliche Informationen, handeln entsprechend optimistisch oder pessimistisch unabhängig davon, ob ihre Einschätzungen objektiv richtig sind.

Diese psychologischen Reaktionen im Zusammenspiel mit Informationsasymmetrien machen es sehr schwer, Aktienbewegungen langfristig exakt vorherzusagen. Denn Preise spiegeln nicht nur fundamentale Daten wider, sondern auch die subjektiven Erwartungen und Unsicherheiten der Marktteilnehmer. So entsteht eine Dynamik, bei der Kursbewegungen weniger vorhersehbar werden, weil unterschiedliche Informationsstände und Wahrnehmungen der Zukunft in den Preis einfließen – oft bevor verlässliche Fakten vorliegen.


\subsubsection{Overfitting und die Illusion von Vorhersagbarkeit}
\paragraph{Vorhersageillusion}
In der Finanzforschung gibt es die Idee, dass Aktienrenditen für Investoren vorhersagbar sein könnten, zum Beispiel basierend auf Kennzahlen wie dem Dividend Preis Verhältnis oder der Entwicklung von Fundamentaldaten. Einige Studien haben gezeigt, dass es statistisch so aussieht, als ob man mit Hilfe bestimmter Modelle billig kaufen und teuer verkaufen könnte. Das führt viele dazu zu glauben, dass Aktienbewegungen prognostizierbar sind. Doch diese Ergebnisse können eine Täuschung sein, die durch Overfitting entsteht, also durch das Überanpassen eines Modells an die historischen Daten, ohne dass es echte Vorhersagekraft hat. In der Studie Predictable Stock Returns: Reality or Statistical Illusion? von Charles R. Nelson und Myung J. Kim wird genau dieses Problem untersucht.
\paragraph{Überanpassung}
Overfitting bedeutet, dass ein Modell nicht nur die wirklich relevanten Zusammenhänge erkennt, sondern auch Zufälligkeiten und sogenanntes Rauschen in den Daten lernt. Das Modell passt sich dann zu stark an die konkreten Werte der Vergangenheit an. Wenn man ein Modell, das perfekt zur Vergangenheit passt, für die Zukunft nutzt, funktioniert es oft weniger gut, weil die gelernten Muster keine echten Gesetzmäßigkeiten widerspiegeln, sondern nur zufällige Schwankungen sind. Nelson und Kim zeigen in ihrer Untersuchung, dass scheinbare Vorhersagemuster allein dadurch entstehen können, dass bestimmte statistische Effekte in kleinen Datensätzen auftreten.
\paragraph{Verzerrung}
Zum Beispiel können Regressionsanalysen mit vergangenen Werten verzerrte Ergebnisse liefern, wenn nur wenige Beobachtungen vorliegen. Dadurch entsteht der Eindruck, dass Renditen vorhersehbar sind, obwohl es in Wirklichkeit keinen stabilen Zusammenhang gibt. Ein weiteres Problem ist die positive Autokorrelation bei mehrjährigen Renditen. Das bedeutet, dass aufeinanderfolgende Zeiträume sich statistisch ähneln können, ohne dass ein echter wirtschaftlicher Zusammenhang besteht. Solche Effekte können dazu führen, dass Ergebnisse als statistisch bedeutsam erscheinen, obwohl sie nur durch Zufall entstanden sind.
\paragraph{Simulation}
Um zu überprüfen, ob diese scheinbaren Muster echte Vorhersagekraft besitzen, führten die Autoren Simulationen mit künstlich erzeugten Daten durch, bei denen Renditen eigentlich nicht vorhersagbar waren. Trotzdem zeigten die statistischen Tests ähnliche Ergebnisse wie in echten Finanzdaten. Das deutet darauf hin, dass viele angebliche Vorhersagemodelle in Wirklichkeit nur statistische Zufälle widerspiegeln.
\paragraph{Prognosekritik}
Für die Bewertung der Prognostizierbarkeit von Aktienbewegungen bedeutet das, dass man sehr vorsichtig sein muss. Nur weil ein Modell in der Vergangenheit gut funktioniert hat, heißt das nicht automatisch, dass es auch in der Zukunft erfolgreich sein wird. Oft entsteht nur die Illusion von Vorhersagbarkeit, weil Modelle zu stark an alte Daten angepasst wurden. Dieses Problem zeigt, wie schwierig es ist, verlässliche Prognosen für Aktienkurse zu erstellen, und warum viele Experten daran zweifeln, dass langfristige Kursbewegungen wirklich zuverlässig vorhergesagt werden können.

\subsection{Warum Forecasting schlecht funktioniert}
\subsubsection{Dominanz von Rauschen und Random-Walk-Charakter von Märkten}

\subsubsection{Nicht stationarität und strukturelle Veränderungen}

\subsubsection{Externe Einflüsse und unvorhersehbare Ereignisse}

\subsection{Interpretation und Einordnung der Ergebnisse}
\subsubsection{Statistische versus ökonomische Signifikanz}

\subsubsection{Grenzen von Performance-Metriken}

\subsubsection{Zusammenhang zwischen ML-Ergebnissen und Markteffizienz}
